\section{Proofs}\label{app:proofs}

Throughout, $[x]^+ = \max\{0,x\}$, payoffs are normalized as in \cref{sec:setup}, and equilibria respect the tie-breaking rules stated there.

\subsection*{Proof of \Cref{lem:credible}}

Stage (4) decisions are taken after Europe's alignment is fixed and after transfers have been paid, so each decision maximizes the decision maker's ex post payoff; by the tie-breaking rule, a power that is indifferent carries out the threat.
(i) In a crisis with cost $\xi$, defending yields the United States $s_a - \xi$ and not defending yields zero, so it defends iff $s_a \geq \xi$, an event of probability $G(s_a)=\pi_a$.
(ii) After exit, denial yields $\nu - k_X$ and avoids the audience cost $m_{ai}$ attached to the public denial rule; supplying yields zero. Denial occurs iff $\nu + m_{ai} \geq k_X$. After Atlantic alignment no commitment is at stake and denial occurs iff $\nu \geq k_A$. Because $k_A>k_X$ and $m_{ai}\geq 0$, denial after $A$ implies denial after $X$, so the conditional part of the instrument is $\chi_{ai} = \mathbf{1}\{\nu+m_{ai}\geq k_X\}\mathbf{1}\{\nu<k_A\}$.
(iii) After exit, penalizing ASML yields $\gamma - \varphi V$ and failing to do so costs $m_f$; the United States penalizes iff $\varphi V \leq \gamma + m_f$. After Atlantic alignment, penalizing yields $-\varphi V<0$ (since $\varphi>0$) and no commitment is at stake.
(iv) After $A$, punishment yields China $g_r$ and avoids audience cost $m_r$; it punishes iff $g_r + m_r\geq 0$. After $X$, punishment yields a negative payoff by assumption.
Europe's expected security loss after $a$ is $(1-\pi_a)S$, so exit adds $(\pi_A - \pi_X)S = \kappa S$. The other terms follow directly, and terms common to both alignments (for instance, AI denial when $\nu\geq k_A$) cancel. \qed

\subsection*{Proof of \Cref{prop:main}}

Fix the grid mesh $\varepsilon>0$. By \cref{lem:credible} and \cref{eq:uE}, Europe chooses $A$ iff $\tau_U - \chi_r r \geq \lambda_C\tau_C - c$.

\emph{China's best response.} Given $\tau_U$, if $\tau_U - \chi_r r + c < 0$, Europe exits even when China offers nothing. Otherwise China wins only with an offer satisfying $\lambda_C \tau_C > \tau_U - \chi_r r + c$; let $\underline\tau_C(\tau_U)$ be the smallest grid point with this property, so that $\underline\tau_C(\tau_U) \in \big((\tau_U - \chi_r r + c)/\lambda_C,\; (\tau_U - \chi_r r + c)/\lambda_C + \varepsilon\big]$. Winning yields $(1-\theta)W - \underline\tau_C(\tau_U)$ and requires $\underline\tau_C(\tau_U)\leq B_C$; losing yields zero; and China does not bid at zero profit. Either way, Europe exits whenever $\tau_U < Q - \lambda_C\varepsilon$ and stays whenever $\tau_U \geq Q$.

\emph{The US offer.} Any grid point $\tau_U \geq [Q]^+$ keeps Europe aligned. The United States' cheapest such offer, $\tau_U^\varepsilon$, satisfies $[Q]^+ - \lambda_C \varepsilon \leq \tau_U^\varepsilon \leq [Q]^+ + \varepsilon$; the lower bound arises because offers in $[Q - \lambda_C\varepsilon, Q)$ may also deter China, depending on where the grid falls. Keeping Europe aligned at this price is strictly better for the United States than losing $V$ iff $\tau_U^\varepsilon < V$, and feasible iff $\tau_U^\varepsilon \leq B_U$. Under $Q<V$ and $Q \leq B_U$ both hold for all small $\varepsilon$, because $\tau_U^\varepsilon \leq [Q]^+ + \varepsilon < V$ and the deterrence threshold is at most $Q \leq B_U$, a grid point. Hence every equilibrium has Europe choosing $A$ and $\tau_U^\varepsilon \to [Q]^+$. Payoffs follow from \cref{eq:uE}. If $Q > \min\{V,B_U\}$, for small $\varepsilon$ the United States either cannot or will not offer $Q$, and China wins. At $Q=V$ the United States is indifferent in the limit, and equilibria with and without exit both exist. \qed

\subsection*{Proof of \Cref{prop:protocol}}

Recall $\bar L = \min\{V, [\lambda_C B_C - c + \chi_r r]^+\}$. Under the stated assumptions, $[Q]^+ \leq \bar L \leq V$ and $L_0 \leq [Q]^+$.

(a) This is \cref{prop:main}.

(b) In any Nash equilibrium in which Europe chooses $A$, the US offer must satisfy $\tau_U \geq [Q]^+$; otherwise China could deviate to a profitable winning offer. The US offer must also equal $[\lambda_C\tau_C - c + \chi_r r]^+$ for China's actual offer $\tau_C$; otherwise the United States could lower its offer and still win. Any $\tau_C \leq B_C$ with $\lambda_C\tau_C - c + \chi_r r \in [[Q]^+, \bar L]$ supports such an equilibrium, with the United States offering exactly that amount: China loses, earns zero, and cannot profitably win because $\tau_U \geq [Q]^+$; the United States prefers matching to losing because the required transfer is at most $V$. China's offers above $(1-\theta)W$ would yield a loss if accepted. If China never makes them, $\lambda_C\tau_C - c + \chi_r r \leq Q$, so $L = [Q]^+$. (On a continuum the offer $(1-\theta)W$ is itself weakly dominated by slightly lower offers, which is why we state the restriction directly; on the $\varepsilon$-grid the limit is the same.)

(c) Given $\tau_C$, the United States offers $[\lambda_C\tau_C - c + \chi_r r]^+$ if that is less than $V$, and otherwise lets Europe exit. Because $Q< V$, China cannot win at a price below its valuation, so China earns zero and is indifferent among all losing offers. Every $\tau_C \in [0,B_C]$ with $\lambda_C\tau_C - c + \chi_r r < V$ is part of a subgame-perfect equilibrium, which gives $L \in [L_0, \bar L]$, with $L_0 = [\chi_r r - c]^+$ obtained at $\tau_C=0$. Excluding offers above $(1-\theta)W$ gives $L \in [L_0, [Q]^+]$.

(d) After a US rejection, China can induce exit at a profit iff $Q>0$, by offering slightly more than $(c-\chi_r r)/\lambda_C$ (or nothing, if $\chi_r r > c$). If $Q>0$ the United States accepts any demand up to $V$, so Europe demands $V$. If $Q\leq 0$ a rejection leaves Europe aligned with no transfer, so the United States rejects every positive demand.

(e) Payoffs are per-period flows. The status quo gives the United States the alignment it wants and gives Europe nothing, so absent exit the United States accepts no positive transfer. Exit, available to Europe only right after it rejects a US offer, yields Europe $\omega - c$ per period against $-\chi_r r$ under the status quo. Suppose $Q>0$. In the stationary equilibrium the United States, when it proposes, offers the per-period transfer $Q$, the least that stops Europe from rejecting and exiting. When Europe proposes, the United States accepts a transfer $\tau$ forever iff $-\tau/(1-\delta) \geq -\delta Q/(1-\delta)$, so Europe demands $\delta Q$. Both converge to $Q$ as $\delta\to 1$. If $Q\leq 0$, exit is not a credible response to a rejection and the United States accepts no positive transfer. This is the outside-option principle \citep{shaked1984involuntary, sutton1986noncooperative}. The restriction on when Europe may exit matters: if Europe could also exit right after the United States rejected Europe's own demand, Europe's continuation after that rejection would be $\delta Q < Q$, exit would be a credible response, and demands up to $V$ could be supported.

(f) Measure Europe's payoffs relative to Atlantic alignment without a transfer, so that agreement on $\tau$ gives Europe $\tau$ and the United States $-\tau$, and breakdown into exit gives Europe $Q$ and the United States $-V$. This adapts the breakdown model of \citet{binmore1986nash} to a random proposer. In a stationary equilibrium with immediate agreement, let $v$ be Europe's expected transfer at the start of a round. When the United States proposes, it offers the least Europe accepts, $\tau_U = \mu Q + (1-\mu)v$. When Europe proposes, it demands the most the United States accepts, $\tau_E = \mu V + (1-\mu)v$. Then $v = \beta\tau_E + (1-\beta)\tau_U = \mu[\beta V + (1-\beta)Q] + (1-\mu)v$, so $v = Q + \beta(V-Q)$ for every $\mu\in(0,1]$, and both proposals converge to $v$ as $\mu \to 0$. With transfers restricted to be nonnegative, $L = [Q + \beta(V-Q)]^+$. If each round also carries discounting with $1-\delta$ comparable to $\mu$, the disagreement point moves toward the status quo; with a fixed $\delta<1$ and status-quo payoffs during bargaining, the limit is that of case (e). \qed

\subsection*{Proof of \Cref{prop:statics}}

(a) By \cref{eq:Q}, $S$ enters only through $\kappa S$, and $\kappa = G(s_A) - G(s_X) = 0$ when $s_A = s_X$. Holding $\pi_X$ fixed, $\partial Q/\partial \pi_A = -S< 0$, and $[Q]^+$ is nondecreasing in $Q$.
(b) $\chi_{ai}$ switches from 0 to 1 at $\nu = k_X - m_{ai}$ and from 1 to 0 at $\nu = k_A$, with $k_X - m_{ai} < k_A$. So $L=[Q_0]^+$ outside the window and $[Q_0-\alpha]^+$ inside it, which is strictly lower iff $\alpha>0$ and $Q_0>0$.
(c) $V$ enters $Q$ only through $\chi_f = \mathbf{1}\{\varphi V \leq \gamma + m_f\}$, which is nonincreasing in $V$; when it switches from 1 to 0, $Q$ rises by $f$.
(d) By \cref{prop:main} Europe's payoff is $[Q]^+ - \chi_r r = \max\{-\chi_r r,\, \omega - c\}$.
(e) If $B_U < Q < V$, the last step of the proof of \cref{prop:main} shows that Europe exits. Any enforceable transfer $\tau\in[Q,V]$ would give Europe $\tau - \chi_r r \geq \omega - c$ and the United States $-\tau \geq -V$. \qed

\subsection*{Proof of \Cref{lem:bite,lem:selflimit}}

After a refusal, cutting yields the United States $g_N - [Q_1+\alpha_1]^+ - \delta[Q_2]^+$ and not cutting yields $-m - [Q_1]^+ - \delta[Q_2]^+$; comparing them, with ties resolved in favour of carrying out the threat, gives \cref{eq:credcut}. Europe's payoff falls by $\alpha_1 - ([Q_1+\alpha_1]^+ - [Q_1]^+)$. If $Q_1\geq 0$ this is zero; if $-\alpha_1 < Q_1 < 0$ it is $-Q_1$; if $Q_1 \leq -\alpha_1$ it is $\alpha_1$. Hence $b_1 = \min\{\alpha_1, [-Q_1]^+\}$. For \cref{lem:selflimit}, the same comparison in period 2 shows that extending the cut costs the United States $[Q_2+\alpha_2]^+ - [Q_2]^+ - (\nu_2 - k_A)$ net of its direct gain, and that Europe's period-2 net loss is $\min\{\alpha_2,[-Q_2]^+\}$, which is zero when $Q_2\geq 0$. \qed

\subsection*{Proof of \Cref{prop:control}}

Suppose (ii) holds, so the United States can choose $m$ large enough to make the cut credible. Europe's refusal payoff is $\rho_E = [Q_1]^+ - b_1 + \delta[Q_2]^+$, and it accepts iff $T - \sigma - r^c \geq \rho_E$. The cheapest acceptable payment is $T^* = [\rho_E + \sigma + r^c]^+$. The United States issues the ultimatum iff its gain $\Gamma \equiv [Q_1]^+ + \delta[Q_2]^+ - T^*$ is strictly positive.

If $b_1 = 0$, then $T^* \geq [Q_1]^+ + \delta[Q_2]^+$ and $\Gamma \leq 0$. If $b_1>0$, then $Q_1<0$, so $\Gamma = \delta[Q_2]^+ - [\delta[Q_2]^+ - b_1 + \sigma + r^c]^+ = \min\{\delta[Q_2]^+,\, b_1 - \sigma - r^c\}$, which is positive iff $Q_2>0$ and $b_1>\sigma + r^c$. If (ii) fails, refusal costs Europe nothing, $T^* = [Q_1]^+ + \delta[Q_2]^+ + \sigma + r^c$, and $\Gamma \leq 0$. Europe's payoff when the ultimatum is issued is $T^* - \sigma - r^c = \max\{\delta Q_2 - b_1, -\sigma - r^c\}$. \qed

\subsection*{Proof of \Cref{cor:windows}}

Call a date $t$ \emph{eligible} if the cut is credible at $t$, $b_t > \sigma + r^c$, and $Q_{t'}>0$ for some $t'>t$. In any subgame starting at a date at or after which no date is eligible, the argument in the proof of \cref{prop:control}, with $\delta[Q_2]^+$ replaced by the discounted sum of later rents, shows that no ultimatum is issued: at each such date either the cut is not credible, or the bite is too small, or there are no later rents to remove, and in each case the gain from issuing is not strictly positive. Let $t^*$ be the last eligible date. In the subgame beginning at $t^*$ with Europe in control, the continuation after a refusal contains strictly positive rents and no further ultimatum, so the United States' gain from issuing is $\min\{\text{discounted future rents},\, b_{t^*} - \sigma - r^c\} > 0$, and it issues. On the equilibrium path an ultimatum is therefore issued at $t^*$ or earlier. Conversely, if no date is eligible, then by backward induction no subgame contains an ultimatum. \qed

\subsection*{Proof of \Cref{prop:breakdown}}

The United States sets $m\geq e$, so it carries out the cut after a refusal. Europe accepts iff $T - \sigma - r^c \geq -b_1 + R$, i.e., iff $\sigma \leq T + \tilde b - R$. Relative to issuing no ultimatum, the United States gains $R - T$ if Europe accepts and loses $e$ if Europe refuses. Write $x = R - T \leq R$, so the acceptance cutoff is $\tilde b - x$. For a cutoff in $(0,\bar\sigma)$ its expected gain is
\begin{equation*}
\Pi(x) = \frac{(\tilde b-x)\,x - (\bar\sigma - \tilde b + x)\,e}{\bar\sigma},
\end{equation*}
which is strictly concave with unconstrained maximizer $x^\circ = (\tilde b-e)/2$, so the cutoff is $(\tilde b+e)/2$. If the cutoff is at least $\bar\sigma$, every type accepts and the gain is $x$, which is increasing in $x$, so the United States never sets a cutoff above $\bar\sigma$. Imposing a cutoff of at most $\bar\sigma$ and $x\leq R$ gives $\hat\sigma = \max\{\tilde b-R, \min\{\bar\sigma, (\tilde b+e)/2\}\}$. At the optimum $x = \tilde b - \hat\sigma$, so the expected gain is $F^*(\tilde b - \hat\sigma) - (1-F^*)e$ with $F^*=\min\{\hat\sigma/\bar\sigma,1\}$, and the United States issues iff this is strictly positive. In the interior this condition is $(\tilde b+e)^2 > 4\bar\sigma e$; if $\hat\sigma=\bar\sigma$ it is $\tilde b>\bar\sigma$; if $\hat\sigma = \tilde b-R<\bar\sigma$ it is $R(\tilde b-R)>(\bar\sigma-\tilde b+R)e$; and if $\tilde b-R\geq\bar\sigma$ it always holds. Conditional on issuance, the comparative statics follow because $\hat\sigma$ is nondecreasing in $\tilde b = b_1 - r^c$ and in $e$, and $P_b = 1 - F^*$. In the corner $\hat\sigma = \tilde b-R < \bar\sigma$, $\hat\sigma$ is strictly decreasing in $R$. \qed

\subsection*{Proof of \Cref{rem:bind}}

If the conditions of \cref{prop:control} hold, Europe's payoff without binding is $\max\{\delta Q_2 - b_1, -\sigma - r^c\}$. With binding, Europe has no future exit premium. By (iii) Europe would accept a demand for control without payment, so the United States gains nothing from issuing one and, under the tie-breaking assumption of \cref{prop:control}, does not issue it. Europe's payoff is $-\sigma^b - r^b$ (period-1 rents are zero because $Q_1<0$). Binding is strictly better iff $-\sigma^b - r^b$ exceeds both terms. \qed

\subsection*{Proof of \Cref{prop:nested}}

(a) If $y < f^F$ the firm refuses to ship after a licence, and exit occurs only if the state compels it, which costs $\zeta$; substituting into \cref{prop:main} gives the result.
(b) China needs only $h$'s consent, so the argument of \cref{prop:main} applies to $h$'s payoffs. A benefit $x_h$ that $h$ receives under $A$ and credibly loses under exit raises $c_h$ by $x_h$, so the payment falls from $[Q_h]^+$ to $[Q_h - x_h]^+$.
(c) With no rival bidder, a countermeasure is adopted iff the members with $q_i>\tau_i$ form a qualified majority, which the United States prevents most cheaply by paying $[q_i]^+$ to every member of some blocking coalition. With a rival bidder moving second, the rival buys a qualified majority $J$ at cost $\sum_{i\in J}[\tau_i - q_i]^+$ whenever that is less than $v_R$, which gives the program in the proposition. Because $[\cdot]^+$ is convex, the program is not linear when some $q_i>0$; it is solved exactly by choosing the set of members with $q_i>0$ whom the United States pays at least $q_i$, setting $\tau_i=0$ for the other members with $q_i>0$, and solving one linear program for each such set. Example: three members, adoption by two, $q_i = 0$. Then $\Phi(0)=0$, while $\Phi(v_R)$ requires $\tau_i + \tau_j \geq v_R$ for every pair, so $\Phi(v_R) = 3v_R/2$ with $\tau_i = v_R/2$ for all three. \qed

\subsection*{Proof of \Cref{prop:holdup}}

(a) Let $p(t)$ be Europe's expected share at date $t$ before the proposer is drawn. A proposer offers the responder exactly its continuation value, and delay is weakly costly to both. Hence
\begin{equation*}
p(t) = \psi \max\{p(t+\Delta) - \hat c_E(t)\Delta,\, 0\} + (1-\psi)\min\{1,\, p(t+\Delta) + \hat c_U(t)\Delta\}.
\end{equation*}
After both buffers are exhausted the costs are constant. In the interior $p(t) - p(t+\Delta) = [(1-\psi)\hat c_U - \psi\hat c_E]\Delta$, so in an infinite horizon the unique solution in $[0,1]$ is $1-\psi\hat c_E\Delta/(1-\psi)$ if $(1-\psi)\hat c_U > \psi\hat c_E$ and $(1-\psi)\hat c_U\Delta/\psi$ otherwise (any other path leaves $[0,1]$ in finite time); these converge to 1 and 0 as $\Delta\to 0$ \citep[cf. the fixed-cost case of][]{rubinstein1982perfect}. Between the two buffer dates only one side bears costs. If $D>0$, Europe alone bears costs on $[t_E,t_U)$, and $p$ falls, moving backward in time from $t_U$, at rate $\psi\hat c_E$, truncated at zero. If $D<0$, the United States alone bears costs on $[t_U,t_E)$, and $p$ rises, moving backward in time from $t_E$, at rate $(1-\psi)\hat c_U$, truncated at one. Before the earlier buffer date nobody bears costs and $p$ is constant. Taking $\Delta\to 0$ gives the stated shares.
(b) Without interruption the United States bears no cost and Europe bears $c_E$ from date zero, so $p^*=0$. With interruption in every round, part (a) with $t_E=0$ gives $p^* = [1 - \psi t_U (c_E+\ell)/z_E]^+ = [1-t_U/\bar t_E]^+$ if Europe is the less vulnerable side post-buffer, and $p^*=0$ otherwise. Because Europe sets interruption at the start of each round, before offers, it interrupts in a round iff doing so raises its expected share from that round on; by backward induction on the remaining US buffer this holds in every round iff it holds at date zero. \qed

\subsection*{Proof of \Cref{prop:hedge}}

After a hedge, the United States compares withholding $j$, which yields $g_j(H)$ but removes $j$ from Europe's exit costs, so that the subsequent contest costs it $[Q + d_j]^+$, with not withholding, which costs $m_j$ if it had committed to punish hedging and leaves the contest at $[Q]^+$. Under $Q + d_j < V$ and $Q+d_j \leq B_U$, \cref{prop:main} applies to the subsequent contest in both cases. Jurisdiction over hedged flows does not enter Europe's exit costs, so enforcing it leaves the contest at $[Q]^+$. \qed

\section{Numerical verification}\label{app:numerics}

We checked every result by brute force in Python; the scripts are in the \texttt{code/checks} folder of the project repository, and none breaks ties in favour of a stated result except as the rules in \cref{sec:setup} require. \Cref{lem:credible} was checked on 200{,}000 draws with stage-4 best-response sets computed without tie-breaking; all generic draws matched. \Cref{prop:main} was checked by grid backward induction that keeps every US-optimal offer, on 3{,}000 random draws with $\lambda_C<1$, capped Chinese offers, and $\theta=1$ allowed; all draws with $Q<\min\{V,B_U\}-\varepsilon$ and all draws in the exit region matched. Cases (a)--(d) of \cref{prop:protocol} were checked by enumerating all pure equilibria on a grid (400 draws), and cases (e) and (f) by exact stationary and finite-horizon computations. \Cref{prop:control} was checked on 200{,}000 draws in which the full set of equilibrium outcomes was computed with every tie admissible; outside the knife-edge where condition (ii) holds with equality, the payment and Europe's payoff matched, and the ``if and only if'' holds under the stated rule that the United States does not issue an ultimatum when indifferent. \Cref{cor:windows} was checked on 100{,}000 draws with two to six periods under the stated rules. \Cref{prop:breakdown} was checked against direct numerical optimization of the US offer with the execution cost $e$ included (40{,}000 draws). \Cref{prop:holdup} was checked by discrete-time backward induction with $\Delta\in\{0.01,0.002\}$ and both terminal conditions; deviations from the closed form are $O(\Delta)$. For part (b) the check treats interruption as chosen for each round before offers are made; under a timing in which Europe can resume servicing after any rejection, interruption is not played. \Cref{prop:nested}(c) was checked against brute-force rival bribing in 3{,}000 small voting games, and the Council arithmetic against the population figures in \citet{council2025population}.
